\documentclass[10pt]{article}
\usepackage[polish]{babel}
\usepackage[utf8]{inputenc}
\usepackage[T1]{fontenc}
\usepackage{amsmath}
\usepackage{amsfonts}
\usepackage{amssymb}
\usepackage[version=4]{mhchem}
\usepackage{stmaryrd}
\usepackage{bbold}
\usepackage{hyperref}
\hypersetup{colorlinks=true, linkcolor=blue, filecolor=magenta, urlcolor=cyan,}
\urlstyle{same}

\title{Construction of APN permutations via Walsh zero spaces }

\author{Benjamin Chase\\
Petr Lisoněk*\\
Department of Mathematics\\
Simon Fraser University\\
Burnaby, BC\\
Canada V5A 1S6\\
plisonek@sfu.ca}
\date{}


%New command to display footnote whose markers will always be hidden
\let\svthefootnote\thefootnote
\newcommand\blfootnotetext[1]{%
  \let\thefootnote\relax\footnote{#1}%
  \addtocounter{footnote}{-1}%
  \let\thefootnote\svthefootnote%
}

%Overriding the \footnotetext command to hide the marker if its value is `0`
\let\svfootnotetext\footnotetext
\renewcommand\footnotetext[2][?]{%
  \if\relax#1\relax%
    \ifnum\value{footnote}=0\blfootnotetext{#2}\else\svfootnotetext{#2}\fi%
  \else%
    \if?#1\ifnum\value{footnote}=0\blfootnotetext{#2}\else\svfootnotetext{#2}\fi%
    \else\svfootnotetext[#1]{#2}\fi%
  \fi
}

\begin{document}
\maketitle


\begin{abstract}
A Walsh zero space (WZ space) for $f: \mathbb{F}_{2^{n}} \rightarrow \mathbb{F}_{2^{n}}$ is an $n$ dimensional vector subspace of $\mathbb{F}_{2^{n}} \times \mathbb{F}_{2^{n}}$ whose all nonzero elements are Walsh zeros of $f$. We provide several theoretical and computerfree constructions of WZ spaces for Gold APN functions $f(x)=x^{2^{i}+1}$ on $\mathbb{F}_{2^{n}}$ where $n$ is odd and $\operatorname{gcd}(i, n)=1$. We also provide several constructions of trivially intersecting pairs of such spaces. We illustrate applications of our constructions that include constructing APN permutations that are CCZ equivalent to $f$ but not extended affine equivalent to $f$ or its compositional inverse.
\end{abstract}

\section*{1 Introduction}
One of the most important open problems concerning APN permutations is their existence in even dimensions. Dillon et al. 5 constructed an APN permutation of $\mathbb{F}_{2^{6}}$ in 2009 but since then, no new APN permutations in even dimension have been found. The question of existence of APN permutations in even dimension greater than six is known as the Big APN Problem.

For odd dimensions, there are six known infinite families of monomial APN permutations, see Table 1 of [6]. An important open problem is the

\footnotetext{*Research of both authors was supported in part by the Natural Sciences and Engineering Research Council of Canada (NSERC).
}
conjecture by Dobbertin [13] that any monomial APN function is CCZequivalent to a permutation belonging to one of these known families.

Currently, all known APN permutations in odd dimensions, up to CCZ equivalence, belong to a monomial family or one of the following families: (1) An infinite family of quadratic APN permutations in odd dimension found in 2008 by Beierle, Carlet, and Leander [8]. (2) Two quadratic APN permutations in dimension 9 found in 2020 by Beierle and Leander [2.

As such, it is certainly still desirable to find new constructions of APN permutations in odd dimensions as well. This is one of the main goals of this paper. The other goal are theoretical descriptions of Walsh zero (WZ) spaces of functions. Some WZ spaces of Gold functions in odd dimensions are constructed computer-free, inspired by concrete examples in low dimensions. WZ spaces already occur implicitly in the work of Dillon et al. [5]; recently their importance and applications were recognized more explicitly, see 9 and references therein.

\section*{2 Preliminaries}
Let $\mathbb{F}_{2^{n}}$ denote the finite field with $2^{n}$ elements and let $\mathbb{F}_{2^{n}}^{*}$ denote the set of its nonzero elements. A function $f: \mathbb{F}_{2^{n}} \rightarrow \mathbb{F}_{2^{n}}$ is almost perfect nonlinear (APN) if for all $a, b \in \mathbb{F}_{2^{n}}, a \neq 0$, the equation $f(x+a)-f(x)=b$ has at most two solutions $x \in \mathbb{F}_{2^{n}}$. Without loss of generality, we can normalize any APN function such that $f(0)=0$, and we will assume this throughout.

APN functions, and more generally functions with low differential uniformity, have been extensively studied due to their importance in the design of S-boxes of block ciphers in cryptography, where they offer the best possible protection against differential cryptanalysis. In some block cipher designs, such as substitution-permutation networks (SPN) [17, Chapter 4], it is required that S -boxes are invertible mappings. Of special interest are therefore APN functions which are invertible, that is, they are permutations of $\mathbb{F}_{2^{n}}$. Constructing new APN permutations of $\mathbb{F}_{2^{n}}$ is one of the objectives of our work.

Let $\operatorname{Tr}_{m}^{n}$ denote the trace function from $\mathbb{F}_{2^{n}}$ to $\mathbb{F}_{2^{m}}$, and let $\operatorname{Tr}$ denote the absolute trace function from $\mathbb{F}_{2^{n}}$ to $\mathbb{F}_{2}$. Let $f$ be a function from $\mathbb{F}_{2^{n}}$ to $\mathbb{F}_{2^{n}}$. For $(a, b) \in \mathbb{F}_{2^{n}} \times \mathbb{F}_{2^{n}}$ we define the Walsh transform of $f$ at ( $a, b$ ) as $\mathcal{W}_{f}(a, b)=\sum_{x \in \mathbb{F}_{2} n}(-1)^{\operatorname{Tr}(a x+b f(x))}$. We say that $(a, b)$ is a Walsh zero of $f$ if $\mathcal{W}_{f}(a, b)=0$.

Definition 2.1. Let $f$ be a function from $\mathbb{F}_{2^{n}}$ to $\mathbb{F}_{2^{n}}$. Suppose that $S$ is an $\mathbb{F}_{2}$-linear subspace of $\mathbb{F}_{2^{n}} \times \mathbb{F}_{2^{n}}$ such that $\operatorname{dim}_{\mathbb{F}_{2}} S=n$ and each element of\\
$S$ except $(0,0)$ is a Walsh zero of $f$. We say that $S$ is a $W Z$ space of $f$.\\
We say that two WZ spaces $S, T$ of the same function intersect trivially if $S \cap T=\{(0,0)\}$.

We will need some background from coding theory. The Hamming weight wt $(x)$ of a vector $x \in \mathbb{F}_{2}^{n}$ is the number of nonzero coordinates in $x$. A binary linear code $C$ with parameters $[n, k, d]$ is an $\mathbb{F}_{2}$-linear subspace of $\mathbb{F}_{2}^{n}$ where $\operatorname{dim}_{\mathbb{F}_{2}} C=k$ and $d$ is the minimum Hamming weight of nonzero codewords in $C$. We say that $d$ is the distance of $C$. A generator matrix for $C$ is a $k \times n$ matrix whose rows form a basis for $C$. Permuting the coordinates of a code $C$ will produce a code $C^{\prime}$ with identical parameters. We say that such codes $C$ and $C^{\prime}$ are equivalent.

For a positive integer $r$, the binary simplex code $\mathcal{S}_{r}$ is a binary linear code with parameters $\left[2^{r}-1, r, 2^{r-1}\right]$. It is well known that this definition is sound, that is, the parameters determine the simplex code uniquely up to equivalence. Moreover for every nonzero $x \in S_{r}$ the Hamming weight of $x$ is $2^{r-1}$. A generator matrix for the code $S_{r}$ is an $r \times\left(2^{r}-1\right)$ matrix whose columns consist of all distinct nonzero vectors of $\mathbb{F}_{2}^{r}$.

The CCZ equivalence of functions was introduced by Carlet, Charpin and Zinoviev in 12. It has many important features, in particular it preserves the APN property. Dillon et al. introduced in 5 a method that, assuming certain conditions are satisfied, constructs a permutation that is CCZ equivalent to a given function. In the following proposition we present this method in a different but equivalent form, using the concept of WZ spaces. We also include a proof of the proposition, which is contained only implicitly in [5], because it allows one to explicitly construct a permutation CCZ equivalent to the given function.

Proposition 2.2. Let $f$ be a function from $\mathbb{F}_{2^{n}}$ to $\mathbb{F}_{2^{n}}$. If there exist two $W Z$ spaces of $f$ that intersect trivially, then $f$ is $C C Z$-equivalent to a permutation of $\mathbb{F}_{2^{n}}$.

Proof. Without loss of generality we can assume that $f(0)=0$, since this can be achieved by an affine transformation which is a special kind of CCZ equivalence.

Let $\left\{\alpha_{1}, \ldots, \alpha_{n}\right\}$ and $\left\{\beta_{1}, \ldots, \beta_{n}\right\}$ be two dual bases of $\mathbb{F}_{2^{n}}$ over $\mathbb{F}_{2}$, that is, $\operatorname{Tr}\left(\alpha_{i} \beta_{j}\right)=\delta_{i, j}$ where for $1 \leq i, j \leq n$ we let $\delta_{i, j}=1$ if $i=j$ and $\delta_{i, j}=0$ otherwise. Let the nonzero elements of $\mathbb{F}_{2^{n}}$ be labelled $x_{1}, \ldots, x_{2^{n}-1}$.

Let $G=\binom{G_{1}}{G_{2}}$ be a $(2 n) \times\left(2^{n}-1\right)$ matrix over $\mathbb{F}_{2}$ where $G_{1}$ and $G_{2}$ are $n \times\left(2^{n}-1\right)$ matrices over $\mathbb{F}_{2}$ defined as follows. The entry in row\\
$i$ and column $j$ of $G_{1}$ is $\operatorname{Tr}\left(\alpha_{i} x_{j}\right)$. The entry in row $i$ and column $j$ of $G_{2}$ is $\operatorname{Tr}\left(\alpha_{i} f\left(x_{j}\right)\right)$. Then the $j$-th column of $G$ is of the form $\binom{x_{j}}{f\left(x_{j}\right)}$ where $x_{j}$ and $f\left(x_{j}\right)$ are represented as $n$-dimensional column vectors with respect to the basis $\left\{\beta_{1}, \ldots, \beta_{n}\right\}$ introduced above. Let $C$ be the binary linear code which is the row space of $G$. Each codeword of $C$ is of the form $\operatorname{Tr}\left(r x_{j}+s f\left(x_{j}\right)\right)_{j=1, \ldots, 2^{n}-1}$ where $r, s$ are fixed elements of $\mathbb{F}_{2^{n}}$.

Let $S$ and $T$ be the two given trivially intersecting WZ spaces, and let $B_{1}=\left\{\left(a_{1}, b_{1}\right), \ldots,\left(a_{n}, b_{n}\right)\right\}$ and $B_{2}=\left\{\left(a_{n+1}, b_{n+1}\right), \ldots,\left(a_{2 n}, b_{2 n}\right)\right\}$ be their bases. Note that $B_{1} \cup B_{2}$ is a basis for $\mathbb{F}_{2^{n}} \times \mathbb{F}_{2^{n}}$.

Let $G^{\prime}=\binom{G_{1}^{\prime}}{G_{2}^{\prime}}$ be a $(2 n) \times\left(2^{n}-1\right)$ matrix over $\mathbb{F}_{2}$ where $G_{1}^{\prime}$ and $G_{2}^{\prime}$ are $n \times\left(2^{n}-1\right)$ matrices over $\mathbb{F}_{2}$ defined as follows. The entry in row $i$ and column $j$ of $G_{1}^{\prime}$ is $\operatorname{Tr}\left(a_{i} x_{j}+b_{i} f\left(x_{j}\right)\right)$. The entry in row $i$ and column $j$ of $G_{2}^{\prime}$ is $\operatorname{Tr}\left(a_{n+i} x_{j}+b_{n+i} f\left(x_{j}\right)\right)$. For $i=1,2$ let $C_{i}^{\prime}$ be the row space of $G_{i}^{\prime}$. Since $S$ and $T$ are WZ spaces for $f$, each nonzero codeword of $C_{i}^{\prime}$ has weight $2^{n-1}$ for $i=1,2$. Thus $C_{1}^{\prime}, C_{2}^{\prime}$ are simplex codes $S_{n}$, and $G_{1}^{\prime}$, $G_{2}^{\prime}$ are formed by pairwise distinct nonzero columns. Thus, with respect to the basis $\left\{\beta_{1}, \ldots, \beta_{n}\right\}$, the columns of $G^{\prime}$ can be viewed as $\binom{x}{g(x)}$ where $x$ runs through all nonzero elements of $\mathbb{F}_{2^{n}}$ and $g(x) \in \mathbb{F}_{2^{n}}$. After letting $g(0)=0$ we see that $g$ is a permutation of $\mathbb{F}_{2^{n}}$.

Let $C^{\prime}$ be the rowspace of $G^{\prime}$. Since $B_{1} \cup B_{2}$ is a basis for $\mathbb{F}_{2^{n}} \times \mathbb{F}_{2^{n}}$, each codeword of $C^{\prime}$ is of the form $\operatorname{Tr}\left(r x_{j}+s f\left(x_{j}\right)\right)_{j=1, \ldots, 2^{n}-1}$ where $r, s$ are fixed elements of $\mathbb{F}_{2^{n}}$. Thus the codes $C$ and $C^{\prime}$ are equal, and functions $f$ and $g$ are CCZ equivalent by [4], see also paragraph 4 on page 380 in [11].

\section*{3 Walsh zeros of Gold functions}
In this section we characterize Walsh zeros of Gold APN permutations. That is, we work in odd dimension throughout. We also define and give some results on certain additive subspaces of $\mathbb{F}_{2^{n}}$ which we call compatible subspaces. These spaces appear while studying Walsh zero spaces of Gold APN permutations.

The following proposition was first proved implicitly by Gold in [14.\\
Proposition 3.1. Suppose $n$ is odd and $f: \mathbb{F}_{2^{n}} \rightarrow \mathbb{F}_{2^{n}}$ with $f(x)=x^{2^{i}+1}$ such that $\operatorname{gcd}(i, n)=1$ (so $f$ is a Gold APN function). Then $(a, b) \in \mathbb{F}_{2^{n}} \times \mathbb{F}_{2^{n}}$ is a Walsh zero of $f$ if and only if $\operatorname{Tr}\left(a b^{-\frac{1}{2^{i}+1}}\right)=0$ or $a \neq b=0$.

Proof. First we present a proof given in 16 which covers the case $b=1$. We then extend it to $b \neq 1$. The notation $\chi(y)=(-1)^{\operatorname{Tr}(y)}$ is introduced only to improve readability. Let $f: \mathbb{F}_{2^{n}} \rightarrow \mathbb{F}_{2^{n}}$ be given by $f(y)=y^{2^{i}+1}$ with $\operatorname{gcd}(i, n)=1$. Then for any $c \in \mathbb{F}_{2^{n}}$,

$$
\begin{aligned}
\mathcal{W}_{f}(A, 1) & =\sum_{y \in \mathbb{F}_{2^{n}}}(-1)^{\operatorname{Tr}(A y+f(y))} \\
& =\sum_{z \in \mathbb{F}_{2^{n}}} \chi\left((z+c)^{2^{i}+1}+A(z+c)\right) \\
& =\sum_{z \in \mathbb{F}_{2^{n}}} \chi\left(z^{2^{i}+1}+c z^{2^{i}}+c^{2^{i}} z+c^{2^{i}+1}+A z+A c\right) \\
& =\chi\left(A c+c^{2^{i}+1}\right) \sum_{z \in \mathbb{F}_{2^{n}}} \chi\left(z^{2^{i}+1}+c^{2^{-i}} z+c^{2^{i}} z+A z\right) \\
& =\chi\left(A c+c^{2^{i}+1}\right) \sum_{z \in \mathbb{F}_{2^{n}}} \chi\left(z^{2^{i}+1}+z(L(c)+A)\right)
\end{aligned}
$$

where $L(c)=c^{2^{i}}+c^{2^{-i}}$.\\
It is clear that $L$ is linear and that $\operatorname{Tr}(L(c))=0$ for all $c$. Since the kernel of $L$ is $\{0,1\}$, the image of $L$ contains all elements of $\mathbb{F}_{2^{n}}$ with trace 0 . So if $\operatorname{Tr}(A)=0$ we can chose $c$ such that $A=L(c)$. With this choice of $c$, we have

$$
W_{f}(A, 1)=\chi\left(A c+c^{2^{i}+1}\right) \sum_{z \in \mathbb{F}_{2^{n}}} \chi\left(z^{2^{i}+1}\right)=0
$$

where $\sum_{z \in \mathbb{F}_{2} n} \chi\left(z^{2^{i}+1}\right)=0$ since $z^{2^{i}+1}$ is a permutation of $\mathbb{F}_{2^{n}}$ when $n$ is odd.

If $\operatorname{Tr}(A)=1$, choose $c$ so that $L(c)=A+1$. Then

$$
W_{f}(A, 1)=\chi\left(A c+c^{2^{i}+1}\right) \sum_{z \in \mathbb{F}_{2^{n}}} \chi\left(z^{2^{i}+1}+z\right)=\chi\left(A c+c^{2^{i}+1}\right) W_{f}(1,1)
$$

The following is proved as Theorem 5 of [16],

$$
W_{f}(1,1)= \begin{cases}+2^{(n+1) / 2} & \text { if } n \equiv \pm 1(\bmod 8) \\ -2^{(n+1) / 2} & \text { if } n \equiv \pm 3(\bmod 8)\end{cases}
$$

So $W_{f}(A, 1)$ is nonzero when $\operatorname{Tr}(A)=1$. Therefore $\mathcal{W}_{f}(A, 1)=0$ if and only if $\operatorname{Tr}(A)=0$.

This result is extended to $b \neq 0,1$ with a change of variables. Let $z=b^{\frac{1}{2^{i}+1}} x$. Then

$$
\begin{aligned}
\mathcal{W}_{f}(A, 1) & =\sum_{z \in \mathbb{F}_{2^{n}}} \chi(A z+f(z)) \\
& =\sum_{x \in \mathbb{F}_{2^{n}}} \chi\left(A b^{\frac{1}{2^{i}+1}} x+b f(x)\right) \\
& =\mathcal{W}_{f}(a, b)
\end{aligned}
$$

where $A$ is chosen such that $a=A b^{\frac{1}{2^{i}+1}}$. The condition $\operatorname{Tr}(A)=0$ becomes $\operatorname{Tr}\left(a b^{-\frac{1}{2^{i}+1}}\right)=0$. Finally, if $a \neq b=0$ then $\mathcal{W}_{f}(a, b)=\sum_{x \in \mathbb{F}_{2^{n}}} \chi(a x)=0$.

If it is clear from the context that $f$ is a Gold APN permutation, then we refer to Proposition 3.1 simply as the Walsh zero test.

Suppose $f: \mathbb{F}_{2^{n}} \rightarrow \mathbb{F}_{2^{n}}$. Since $\operatorname{Tr}(a x)$ is a balanced function of $x$, the space

$$
Z_{a 0}=\left\{(a, 0): a \in \mathbb{F}_{2^{n}}\right\}
$$

is a WZ space of $f$. Furthermore, if $f$ is a permutation and $b \neq 0$ then $\operatorname{Tr}(b f(x))$ is a balanced function of $x$, hence and

$$
Z_{0 b}=\left\{(0, b): b \in \mathbb{F}_{2^{n}}\right\}
$$

is also a WZ space of $f$. We call these two spaces trivial WZ spaces.\\[0pt]
WZ spaces have been recently appearing in literature. In Theorem 4 of [9], Canteaut and Perrin prove that the number of EA-classes inside the CCZ-class of $f: \mathbb{F}_{2^{n}} \rightarrow \mathbb{F}_{2^{m}}$ is upper bounded by the number of WZ spaces of $f$. Beierle, Carlet, Leander, and Perrin in 1 have studied two new quadratic APN permutations in dimension 9 using numerical properties of their WZ spaces (thickness and degree spectrum). Furthermore, in 1 the authors show an interesting similarity between their new APN permutations and Gold APN permutations in odd dimension divisible by 3. In the next section we provide constructions of nontrivial WZ spaces for Gold APN permutations in odd dimension which are, as far as we know, the first such theoretical (computer-free) constructions.

From now on let us define $0^{-\frac{1}{2^{i}+1}}=0$ as this will simplify some of the forthcoming constructions and arguments. Let us note that this works correctly for the Walsh zero test for Gold APN permutations, where for a pair ( $a, 0$ ) we get $\operatorname{Tr}\left(a \cdot 0^{-\frac{1}{2^{i}+1}}\right)=\operatorname{Tr}(a \cdot 0)=0$, as desired.

Definition 3.2. Let $n$ be odd and $\operatorname{gcd}(i, n)=1$. Let $S$ be an additive subspace of $\mathbb{F}_{2^{n}}$. We say that $S$ is $i$-compatible if the set $S^{-\frac{1}{2^{i}+1}}=\left\{s^{-\frac{1}{2^{i}+1}}\right.$ : $s \in S\}$ is also an additive subspace of $\mathbb{F}_{2^{n}}$.

For $U \subseteq \mathbb{F}_{2^{n}}$ and $a \in \mathbb{F}_{2^{n}}$ denote $a U=\{a u: u \in U\}$.

\section*{Example 3.3.}
(i) If $n$ is odd and $\operatorname{gcd}(i, n)=1$ then the following subspaces of $\mathbb{F}_{2^{n}}$ are $i$-compatible: $\{0\}, \mathbb{F}_{2}$ and $\mathbb{F}_{2^{n}}$.\\
(ii) If $S$ is an $i$-compatible subspace of $\mathbb{F}_{2^{n}}$, then $\mu S$ is also $i$-compatible for each $\mu \in \mathbb{F}_{2^{n}}$.

Proposition 3.4. Suppose $m$ divides $n$ and $\operatorname{gcd}(i, n)=1$. Then $\mathbb{F}_{2^{m}}$ is $i$-compatible subspace of $\mathbb{F}_{2^{n}}$.

Proof. Let $\xi \in \mathbb{F}_{2^{m}}$. Then $\left(\xi^{-\frac{1}{2^{i}+1}}\right)^{2^{m}}=\left(\xi^{2^{m}}\right)^{-\frac{1}{2^{i}+1}}=\xi^{-\frac{1}{2^{i}+1}}$. So $\xi^{-\frac{1}{2^{i}+1}} \in \mathbb{F}_{2^{m}}$ and it follows that $\left\{\xi^{-\frac{1}{2^{i}+1}}: \xi \in \mathbb{F}_{2^{m}}\right\}=\mathbb{F}_{2^{m}}$.

Moreover, when $n$ is a multiple of 3 , then $\mathbb{F}_{2^{n}}$ contains the subfield $\mathbb{F}_{2^{3}}$ and the following lemma applies.

Lemma 3.5. Suppose $n$ is a multiple of 3 and $\operatorname{gcd}(i, n)=1$. Let $\xi \in \mathbb{F}_{2^{3}} \subset \mathbb{F}_{2^{n}}$. Then $\xi^{-\frac{1}{2^{i}+1}}=\xi^{2^{i}}$.

Proof. For $\xi=0$ the result follows from our definition $0^{-\frac{1}{2^{i}+1}}=0$. Otherwise, after raising both sides of $\xi^{-\frac{1}{2^{i}+1}}=\xi^{2^{i}}$ to power $2^{i}+1$ we face proving the equivalent equation

$$
\xi^{-1}=\xi^{2^{2 i}+2^{i}}
$$

For $i \equiv 1(\bmod 3)$ this follows from

$$
\xi^{2^{2 i}+2^{i}}=\xi^{2^{2}+2^{1}}=\xi^{6}=\xi^{-1}
$$

and for $i \equiv 2(\bmod 3)$ it follows from

$$
\xi^{2^{2 i}+2^{i}}=\xi^{2^{1}+2^{2}}=\xi^{6}=\xi^{-1}
$$

Corollary 3.6. Assume that $n$ is odd and divisible by 3. Let $S$ be the subfield of $\mathbb{F}_{2^{n}}$ isomorphic to $\mathbb{F}_{2^{3}}$. Then each additive subspace of $S$ is $i$-compatible whenever $\operatorname{gcd}(i, n)=1$.

Proof. This follows from Lemma 3.5 since $\xi \mapsto \xi^{-\frac{1}{2^{i}+1}}$ is a linear mapping on $\mathbb{F}_{2^{3}}$.

As any two 2 -dimensional subspaces (hyperplanes) of $\mathbb{F}_{2^{3}}$ can be obtained from each other just by scaling by an element of $\mathbb{F}_{2^{3}}^{*}$, it follows that if 3 divides $n$, then Lemma 3.6 provides $2^{n}-1$ two-dimensional $i$ compatible subspaces of $\mathbb{F}_{2^{n}}$ and $\left(2^{n}-1\right) / 7$ three-dimensional $i$-compatible subspaces of $\mathbb{F}_{2^{n}}$. We have checked by exhaustive search that there are no two-dimensional $i$-compatible subspaces of $\mathbb{F}_{2^{n}}$ other than those described above for odd $n$ and relatively prime $i$ in the range $n \leq 17$.

For $i=1$ it is easy to see that the proof of Lemma 3.6 does not extend to any higher dimension. Indeed suppose that $s \mapsto s^{-1 / 3}$ is linear on $\mathbb{F}_{2^{n}}$ where $n>3$ is odd, then $-1 / 3 \equiv 2^{k}\left(\bmod 2^{n}-1\right)$ for some integer $0 \leq k<n$, and $3 \cdot 2^{k}+1$ is a multiple of $2^{n}-1$. If $n>3$ and $k<n-1$ then $3 \cdot 2^{k}+1$ is too small for this to happen. The only remaining possibility is $3 \cdot 2^{n-1}+1=2^{n}-1$, which can not occur either. Similarly, to extend Lemma 3.6 to higher dimension for $i>1$, we must show that $2^{k}\left(2^{i}+1\right)+1$ is a multiple of $2^{n}-1$ for some integer $0 \leq k<n$. But a parity argument shows that this is also not possible.

In Proposition 4.1 below we will see that $i$-compatible spaces can be used to construct WZ spaces. Thus it is interesting to obtain as many $i$ compatible spaces as we can. This motivates us to present the following open problem:

Problem 3.7. For odd $n$, do there exist $i$-compatible subspaces of $\mathbb{F}_{2^{n}}$ other than those described by Example 3.3 (i,ii), Proposition 3.4 and Corollary 3.6?

\section*{4 Constructions of some WZ spaces for Gold APN permutations}
Theoretical constructions in this section have been partially motivated by computational examples in low odd dimensions obtained using the sboxU software package 15, 3 written by Léo Perrin and Mathias Joly. In particular, we verified using sboxU that the theoretical constructions given in this section, along with the two trivial spaces $Z_{a 0}$ and $Z_{0 b}$, cover all WZ spaces of Gold APN functions in all odd dimensions less than or equal to 9.

Proposition 4.1. Assume that $n$ is odd and $S$ is an $i$-compatible additive subspace of $\mathbb{F}_{2^{n}}$. Let $f: \mathbb{F}_{2^{n}} \rightarrow \mathbb{F}_{2^{n}}, f(x)=x^{2^{i}+1}$, with $\operatorname{gcd}(i, n)=1$. Let

$$
X=\left\{x \in \mathbb{F}_{2^{n}}:\left(\forall a \in S^{-\frac{1}{2^{i}+1}}\right) \operatorname{Tr}(a x)=0\right\}
$$

Then $X \times S$ is a $W Z$ space for $f$.\\
Proof. Let $(0,0) \neq(x, s) \in X \times S$. Then $\operatorname{Tr}\left(x s^{-\frac{1}{2^{i}+1}}\right)=0$ by construction, hence $(x, s)$ is a Walsh zero of $f$. Since both $X$ and $S$ are linear spaces, $X \times S$ is also a linear space. Finally, since $s \mapsto s^{-\frac{1}{2^{i}+1}}$ is a bijection on $\mathbb{F}_{2^{n}}$, we get $\operatorname{dim} S^{-\frac{1}{2^{2}+1}}=\operatorname{dim} S$, and

$$
\operatorname{dim}(X \times S)=\operatorname{dim} X+\operatorname{dim} S=(n-\operatorname{dim} S)+\operatorname{dim} S=n
$$

Proposition 4.2. Assume that $n=3 k$ where $k$ is odd. Let $f: \mathbb{F}_{2^{n}} \rightarrow \mathbb{F}_{2^{n}}$, $f(x)=x^{2^{i}+1}$, with $\operatorname{gcd}(i, n)=1$. Suppose $\xi$ is a fixed element of $\mathbb{F}_{2^{3}} \subset \mathbb{F}_{2^{n}}$ and $\mu$ is a fixed element of $\mathbb{F}_{2^{n}}^{*}$. Then

$$
Z=\left\{\left(x, \mu^{-\left(2^{i}+1\right)}\left(\xi \operatorname{Tr}(\mu x)+\operatorname{Tr}\left(\xi^{2^{i}} \mu x\right)\right)\right): x \in \mathbb{F}_{2^{n}}\right\}
$$

is a $W Z$ space of $f$.\\
Proof. The additive closure of $Z$ follows from the fact that

$$
g(x)=\mu^{-\left(2^{i}+1\right)}\left(\xi \operatorname{Tr}(\mu x)+\operatorname{Tr}\left(\xi^{2^{i}} \mu x\right)\right)
$$

is an $\mathbb{F}_{2}$-linear function.\\
We now prove that each element $(x, g(x)) \in Z \backslash\{(0,0)\}$ is a Walsh zero of $f$. We only need to address the cases when $g(x) \neq 0$. Note that $\mu^{2^{i}+1} g(x) \in \mathbb{F}_{2^{3}}$. Using Lemma 3.5 we get

$$
\begin{aligned}
\operatorname{Tr}\left(x g(x)^{-\frac{1}{2^{i}+1}}\right) & =\operatorname{Tr}\left(x \mu\left(\mu^{2^{i}+1} g(x)\right)^{-\frac{1}{2^{i}+1}}\right) \\
& =\operatorname{Tr}\left(x \mu\left(\xi \operatorname{Tr}(\mu x)+\operatorname{Tr}\left(\xi^{2^{i}} \mu x\right)\right)^{2^{i}}\right) \\
& =\operatorname{Tr}\left(x \mu \xi^{2^{i}} \operatorname{Tr}(\mu x)+x \mu \operatorname{Tr}\left(\xi^{2^{i}} \mu x\right)\right) \\
& =\operatorname{Tr}\left(x \mu \xi^{2^{i}} \operatorname{Tr}(\mu x)\right)+\operatorname{Tr}\left(x \mu \operatorname{Tr}\left(\xi^{2^{i}} \mu x\right)\right) \\
& =0
\end{aligned}
$$

because $\operatorname{Tr}\left(x \mu \xi^{2^{i}} \operatorname{Tr}(\mu x)\right)=\operatorname{Tr}\left(x \mu \xi^{2^{i}}\right) \operatorname{Tr}(\mu x)=\operatorname{Tr}\left(x \mu \operatorname{Tr}\left(\xi^{2^{i}} \mu x\right)\right)$. By considering the first components of the elements of $Z$ it is clear that $\operatorname{dim} Z= n$.

While Proposition 4.2 holds for each $\xi \in \mathbb{F}_{2^{3}}$, one obtains interesting results only when $\xi$ is a primitive element of $\mathbb{F}_{2^{3}}$. If $\xi=0,1$ then $Z$ is the trivial WZ space $Z_{a 0}$.

Proposition 4.3. Suppose $n$ is odd and $m$ divides $n$. Let $f: \mathbb{F}_{2^{n}} \rightarrow \mathbb{F}_{2^{n}}$, $f(x)=x^{2^{i}+1}$, with $\operatorname{gcd}(i, n)=1$. Suppose $\mu$ is a fixed element of $\mathbb{F}_{2^{n}}^{*}$. Then

$$
Z=\left\{\left(\mu a, \mu^{2^{i}+1} b\right): a \in \mathbb{F}_{2^{n}}, b \in \mathbb{F}_{2^{m}} \text { and } \operatorname{Tr}_{m}^{n}(a)=b+b^{\frac{1}{2^{i}}}\right\}
$$

is a $W Z$ space of $f$.\\
Proof. First we show that each element $\left(\mu a, \mu^{2^{i}+1} b\right)$ of $Z$ except for $(0,0)$ is a Walsh zero of $f$. Note that since $m$ is odd, $b^{-\frac{1}{2^{i}+1}} \in \mathbb{F}_{2^{m}}$. We have

$$
\begin{aligned}
\operatorname{Tr}_{1}^{n}\left(\mu a\left(\mu^{2^{i}+1} b\right)^{-\frac{1}{2^{i}+1}}\right) & =\operatorname{Tr}_{1}^{n}\left(a b^{-\frac{1}{2^{i}+1}}\right) \\
& =\operatorname{Tr}_{1}^{m}\left(\operatorname{Tr}_{m}^{n}\left(a b^{-\frac{1}{2^{i}+1}}\right)\right) \\
& =\operatorname{Tr}_{1}^{m}\left(b^{-\frac{1}{2^{i}+1}} \operatorname{Tr}_{m}^{n}(a)\right) \\
& =\operatorname{Tr}_{1}^{m}\left(b^{-\frac{1}{2^{i}+1}}\left(b+b^{\frac{1}{2^{i}}}\right)\right) \\
& =\operatorname{Tr}_{1}^{m}\left(b^{\frac{2^{i}}{2^{i}+1}}+b^{\frac{1}{2^{i}\left(2^{i}+1\right)}}\right) \\
& =\operatorname{Tr}_{1}^{m}\left(b^{\frac{2^{i}}{2^{i}+1}}\right)+\operatorname{Tr}_{1}^{m}\left(b^{\frac{2^{2 i}}{2^{i}\left(2^{i}+1\right)}}\right) \\
& =0
\end{aligned}
$$

hence each nonzero element of $Z$ is a Walsh zero of $f$.\\
Let $\left(\mu a_{1}, \mu^{2^{i}+1} b_{1}\right),\left(\mu a_{2}, \mu^{2^{i}+1} b_{2}\right) \in Z$. Then

$$
\left(\mu a_{1}, \mu^{2^{i}+1} b_{1}\right)+\left(\mu a_{2}, \mu^{2^{i}+1} b_{2}\right)=\left(\mu\left(a_{1}+a_{2}\right), \mu^{2^{i}+1}\left(b_{1}+b_{2}\right)\right) \in Z
$$

since $b_{1}+b_{2} \in \mathbb{F}_{2^{m}}$ and

$$
\begin{aligned}
\operatorname{Tr}_{m}^{n}\left(a_{1}+a_{2}\right) & =\operatorname{Tr}_{m}^{n}\left(a_{1}\right)+\operatorname{Tr}_{m}^{n}\left(a_{2}\right) \\
& =b_{1}+b_{1}^{\frac{1}{2^{i}}}+b_{2}+b_{2}^{\frac{1}{2^{i}}} \\
& =\left(b_{1}+b_{2}\right)+\left(b_{1}+b_{2}\right)^{\frac{1}{2^{i}}}
\end{aligned}
$$

Therefore $Z$ is additively closed.\\
For each of the $2^{m}$ choices for the second component $\mu^{2^{i}+1} b$ of an element of $Z$, there are $2^{n-m}$ choices for the first component $\mu a$ such that $\operatorname{Tr}_{m}^{n}(a)= b+b^{\frac{1}{2^{i}}}$. It follows that the dimension of $Z$ is $n$.

\section*{5 Trivially intersecting pairs of WZ spaces for Gold APN permutations}
In this section we describe several theoretical constructions of trivially intersecting pairs of WZ spaces for Gold APN permutations in odd dimensions. We checked using the sboxU software package that these constructions cover all such pairs in dimensions less than or equal to 9 . Obviously, the two trivial WZ spaces $Z_{a 0}, Z_{0 b}$ intersect trivially.

Proposition 5.1. Let $n=3 k$ where $k$ is odd and $f: \mathbb{F}_{2^{n}} \rightarrow \mathbb{F}_{2^{n}}, f(x)= x^{2^{i}+1}$, with $\operatorname{gcd}(i, n)=1$. Let

$$
Z=\left\{\left(x, \mu^{-\left(2^{i}+1\right)}\left(\xi \operatorname{Tr}(\mu x)+\operatorname{Tr}\left(\xi^{2^{i}} \mu x\right)\right)\right): x \in \mathbb{F}_{2^{n}}\right\}
$$

where $\xi$ is a fixed element of $\mathbb{F}_{2^{3}} \subset \mathbb{F}_{2^{n}}$ and $\mu$ is a fixed element of $\mathbb{F}_{2^{n}}^{*}$. Then $Z$ is a $W Z$ space of $f$ and the pair $\left\{Z_{0 b}, Z\right\}$ intersects trivially.

Proof. It follows from Proposition 4.2 that $Z$ is a WZ space of $f$. If $(x, c) \in Z_{0 b} \cap Z$, then $x=0$ which implies $c=0$. Therefore the pair $\left\{Z_{0 b}, Z\right\}$ intersects trivially.

Note that in the above proposition, if $\xi \in\{0,1\}$ then $Z$ becomes the trivial WZ space $Z_{a 0}$.

Proposition 5.2. Let $n=3 k$ where $k$ is odd and $f: \mathbb{F}_{2^{n}} \rightarrow \mathbb{F}_{2^{n}}, f(x)= x^{2^{i}+1}$, with $\operatorname{gcd}(i, n)=1$. Let

$$
Y=\left\{\left(x, \mu^{-\left(2^{i}+1\right)}\left(\xi \operatorname{Tr}(\mu x)+\operatorname{Tr}\left(\xi^{2^{i}} \mu x\right)\right)\right): x \in \mathbb{F}_{2^{n}}\right\}
$$

where $\xi$ is a fixed element of $\mathbb{F}_{2^{3}} \subset \mathbb{F}_{2^{n}}$ and $\mu$ is a fixed element of $\mathbb{F}_{2^{n}}^{*}$. Let

$$
Z=\left\{\left(\nu\left(b+b^{\frac{1}{2^{i}}}\right), \nu^{2^{i}+1} b\right): b \in \mathbb{F}_{2^{n}}\right\}
$$

with $\nu$ a fixed element of $\mathbb{F}_{2^{n}}^{*}$. Suppose also that $\operatorname{Tr}\left(\left(\xi+\xi^{2^{i}}\right)(\mu \nu)^{-2^{i}}\right)=0$. Then $Y$ and $Z$ are $W Z$ spaces of $f$ and the pair $\{Y, Z\}$ intersects trivially.

Proof. It follows from Proposition 4.2 that $Y$ is a WZ space of $f$ and from Proposition 4.3 with $m=n$ that $Z$ is a WZ space of $f$.

Suppose towards a contradiction that there exists an element, different from $(0,0)$, in $Y \cap Z$. Then for this element, we have $x=\nu\left(b+b^{\frac{1}{2^{i}}}\right)$. It follows that $b \neq 0$ and


\begin{equation*}
(\mu \nu)^{2^{i}+1} b=\xi \operatorname{Tr}\left(\mu \nu\left(b+b^{\frac{1}{2^{i}}}\right)\right)+\operatorname{Tr}\left(\xi^{2^{i}} \mu \nu\left(b+b^{\frac{1}{2^{i}}}\right)\right) . \tag{1}
\end{equation*}


By looking at the right-hand side of (11) we can see that $(\mu \nu)^{2^{i}+1} b \in\{0,1, \xi, \xi+ 1\}$. But since $\mu \neq 0, \nu \neq 0$, and $b \neq 0$, we have $(\mu \nu)^{2^{i}+1} b \in\{1, \xi, \xi+1\}$. We will look at these three cases below.

First assume that $(\mu \nu)^{2^{i}+1} b=1$. Observe that

$$
\begin{aligned}
\operatorname{Tr}\left(\xi^{2^{i}} \mu \nu\left(b+b^{\frac{1}{2^{i}}}\right)\right) & =\operatorname{Tr}\left(\xi^{2^{i}} \mu \nu(\mu \nu)^{-\left(2^{i}+1\right)}+\xi^{2^{i}} \mu \nu(\mu \nu)^{-\frac{2^{i}+1}{2^{i}}}\right) \\
& =\operatorname{Tr}\left(\xi^{2^{i}}(\mu \nu)^{-2^{i}}\right)+\operatorname{Tr}\left(\xi^{2^{i}}(\mu \nu)^{-\frac{1}{2^{i}}}\right) \\
& =\operatorname{Tr}\left(\xi^{2^{i}}(\mu \nu)^{-2^{i}}\right)+\operatorname{Tr}\left(\xi^{2^{3 i}}(\mu \nu)^{-\frac{2^{2 i}}{2^{i}}}\right) \\
& =\operatorname{Tr}\left(\xi^{2^{i}}(\mu \nu)^{-2^{i}}\right)+\operatorname{Tr}\left(\xi^{2^{3 i}}(\mu \nu)^{-2^{i}}\right) \\
& =\operatorname{Tr}\left(\left(\xi^{2^{i}}+\xi\right)(\mu \nu)^{-2^{i}}\right)=0
\end{aligned}
$$

which contradicts (1).\\
Finally assume that $(\mu \nu)^{2^{i}+1} b=z$ and $z \in\{\xi, \xi+1\}$. Then

$$
\begin{aligned}
\operatorname{Tr}\left(\mu \nu\left(b+b^{\frac{1}{2^{i}}}\right)\right) & =\operatorname{Tr}\left(z \mu \nu(\mu \nu)^{-\left(2^{i}+1\right)}+z^{\frac{1}{2^{i}}} \mu \nu(\mu \nu)^{-\frac{2^{i}+1}{2^{i}}}\right) \\
& =\operatorname{Tr}\left(z(\mu \nu)^{-2^{i}}\right)+\operatorname{Tr}\left(z^{\frac{1}{2^{i}}}(\mu \nu)^{-\frac{1}{2^{i}}}\right) \\
& =\operatorname{Tr}\left(z(\mu \nu)^{-2^{i}}\right)+\operatorname{Tr}\left(z^{\frac{2^{2 i}}{2^{i}}}(\mu \nu)^{-\frac{2^{2 i}}{2^{i}}}\right) \\
& =\operatorname{Tr}\left(z(\mu \nu)^{-2^{i}}\right)+\operatorname{Tr}\left(z^{2^{i}}(\mu \nu)^{-2^{i}}\right) \\
& =\operatorname{Tr}\left(\left(\xi+\xi^{2^{i}}\right)(\mu \nu)^{-2^{i}}\right)=0
\end{aligned}
$$

which again contradicts (1).\\
Since all cases are exhausted, the only element in the intersection is (0,0) and the proof is complete.

Proposition 5.3. Let $n=3 k$ where $k$ is odd and $f: \mathbb{F}_{2^{n}} \rightarrow \mathbb{F}_{2^{n}}, f(x)= x^{2^{i}+1}$, with $\operatorname{gcd}(i, n)=1$. Let

$$
Y=\left\{\left(\nu\left(b+b^{\frac{1}{2^{i}}}\right), \nu^{2^{i}+1} b\right): b \in \mathbb{F}_{2^{n}}\right\}
$$

with $\nu$ a fixed element of $\mathbb{F}_{2^{n}}^{*}$. Suppose $\xi$ is a fixed primitive element of $\mathbb{F}_{2^{3}} \subset \mathbb{F}_{2^{n}}$. Let

$$
Z=X \times S
$$

with $S=\operatorname{span}_{\mathbb{F}_{2}}\{\mu, \xi \mu\}$ for some fixed $\mu \in \mathbb{F}_{2^{n}}^{*}$ and $Z$ is constructed by applying Proposition 4.1. Suppose also that $\operatorname{Tr}\left(\left(\xi+\xi^{2^{i}}\right) \mu^{\frac{2^{i}}{2^{i}+1}} \nu^{-2^{i}}\right)=1$. Then $Y$ and $Z$ are $W Z$ spaces of $f$ and the pair $\{Y, Z\}$ intersects trivially.

Proof. It follows from Proposition 4.3 with $m=n$ that $Y$ is a WZ space of $f$.

Note that $S$ is $i$-compatible by Lemma 3.6 and Example 3.3(ii): If $T= \operatorname{span}_{\mathbb{F}_{2}}\{1, \xi\}$ then $S=\mu T$. From Proposition 4.1 we have

$$
X=\left\{x \in \mathbb{F}_{2^{n}}:\left(\forall \alpha \in S^{-\frac{1}{2^{i}+1}}\right) \operatorname{Tr}(\alpha x)=0\right\}
$$

where

$$
S^{-\frac{1}{2^{i}+1}}=\left\{0, \mu^{-\frac{1}{2^{i}+1}},(\xi \mu)^{-\frac{1}{2^{i}+1}},((1+\xi) \mu)^{-\frac{1}{2^{i}+1}}\right\}
$$

Suppose towards a contradiction that $\left(\nu\left(b+b^{\frac{1}{2^{i}}}\right), \nu^{2^{i}+1} b\right)$ is a nonzero element belonging to $Y \cap Z$. Then $\nu^{2^{i}+1} b \in S \backslash\{0\}=\{\mu, \xi \mu,(1+\xi) \mu\}$. We will look at these three cases below.

First assume $\nu^{2^{i}+1} b=\mu$. Take $\alpha=(\xi \mu)^{-\frac{1}{2^{i}+1}} \in S^{-\frac{1}{2^{i}+1}}$. By Lemma 3.5 we have $\alpha=\xi^{2^{i}} \mu^{-\frac{1}{2^{i}+1}}$. Then by the condition defining $X$, we should have $\operatorname{Tr}\left(\alpha \nu\left(b+b^{\frac{1}{2^{i}}}\right)\right)=0$. But

$$
\begin{aligned}
\operatorname{Tr}\left(\alpha \nu\left(b+b^{\frac{1}{2^{i}}}\right)\right) & =\operatorname{Tr}\left(\xi^{2^{i}} \mu^{-\frac{1}{2^{i}+1}} \nu\left(\mu \nu^{-2^{i}-1}+\mu^{\frac{1}{2^{i}}} \nu^{-\frac{1}{2^{i}}-1}\right)\right) \\
& =\operatorname{Tr}\left(\xi^{2^{i}}\left(\mu^{\frac{2^{i}}{2^{2}+1}} \nu^{-2^{i}}+\mu^{\frac{1}{\left(2^{i}+1\right) 2^{i}}} \nu^{-\frac{1}{2^{i}}}\right)\right) \\
& =\operatorname{Tr}\left(\xi^{2^{i}} \mu^{\frac{2^{i}}{2^{i}+1}} \nu^{-2^{i}}\right)+\operatorname{Tr}\left(\xi^{2^{3 i}} \mu^{\frac{2^{2 i}}{\left(2^{i}+1\right) 2^{i}}} \nu^{-\frac{2^{2 i}}{2^{i}}}\right) \\
& =\operatorname{Tr}\left(\xi^{2^{i}} \mu^{\frac{2^{i}}{2^{i}+1}} \nu^{-2^{i}}\right)+\operatorname{Tr}\left(\xi \mu^{\frac{2^{i}}{2^{i}+1}} \nu^{-2^{i}}\right) \\
& =\operatorname{Tr}\left(\left(\xi+\xi^{2^{i}}\right) \mu^{\frac{2^{i}}{2^{i}+1}} \nu^{-2^{i}}\right)=1
\end{aligned}
$$

Finally assume that $\nu^{2^{i}+1} b=z \mu$ and $z \in\{\xi, \xi+1\}$. Take $\alpha=\mu^{-\frac{1}{2^{i}+1}} \in S^{-\frac{1}{2^{i}+1}}$. By the condition defining $X$, we should have $\operatorname{Tr}\left(\alpha \nu\left(b+b^{\frac{1}{2^{i}}}\right)\right)=0$. But

$$
\begin{aligned}
\operatorname{Tr}\left(\alpha \nu\left(b+b^{\frac{1}{2^{i}}}\right)\right) & =\operatorname{Tr}\left(\mu^{-\frac{1}{2^{i}+1}} \nu\left(z \mu \nu^{-2^{i}-1}+z^{\frac{1}{2^{i}}} \mu^{\frac{1}{2^{i}}} \nu^{-\frac{1}{2^{i}}-1}\right)\right) \\
& =\operatorname{Tr}\left(z \mu^{\frac{2^{i}}{2^{i}+1}} \nu^{-2^{i}}+z^{\frac{1}{2^{i}}} \mu^{\frac{1}{\left(2^{i}+1\right) 2^{i}}} \nu^{-\frac{1}{2^{i}}}\right) \\
& =\operatorname{Tr}\left(z \mu^{\frac{2^{i}}{2^{i}+1}} \nu^{-2^{i}}\right)+\operatorname{Tr}\left(z^{\frac{2^{2 i}}{2^{i}}} \mu^{\frac{2^{2 i}}{\left(2^{i}+1\right) 2^{i}}} \nu^{-\frac{2^{2 i}}{2^{i}}}\right) \\
& =\operatorname{Tr}\left(z \mu^{\frac{2^{i}}{2^{i}+1}} \nu^{-2^{i}}\right)+\operatorname{Tr}\left(z^{2^{i}} \mu^{\frac{2^{i}}{2^{i}+1}} \nu^{-2^{i}}\right) \\
& =\operatorname{Tr}\left(\left(\xi+\xi^{2^{i}}\right) \mu^{\frac{2^{i}}{2^{i}+1}} \nu^{-2^{i}}\right)=1
\end{aligned}
$$

Since all cases are exhausted, the only element in the intersection is ( 0,0 ) and we are done.

In Propositions 5.2 and 5.3 it is not difficult to see that if $\xi$ is drawn uniformly at random from $\mathbb{F}_{2^{3}} \backslash\{0,1\}$ and $\mu, \nu$ are drawn uniformly at random from $\mathbb{F}_{2^{n}}^{*}$, then the probability that the trace condition is satisfied is close to $1 / 2$. In other words, we expect that a trivially intersecting pair of WZ spaces is obtained for about every other set of random parameters generated in this way.

\section*{6 Applications}
\subsection*{6.1 Classifying EA classes of functions}
Functions $f$ and $g$ mapping $\mathbb{F}_{2^{n}}$ to $\mathbb{F}_{2^{n}}$ are extended affine equivalent ( $E A$ equivalent) if there exist affine permutations $A_{1}, A_{2}$ of $\mathbb{F}_{2^{n}}$ and an affine function $A_{3}$ such that $A_{1}\left(f\left(A_{2}(x)\right)\right)+A_{3}(x)=g(x)$ for all $x \in \mathbb{F}_{2^{n}}$. It is known that EA equivalent functions are also CCZ equivalent, but partitioning CCZ classes into EA classes is in general a hard problem. This problem was addressed by Canteaut and Perrin 9 by studying the structure of Walsh zeros of functions. WZ spaces play an important role in their investigations.

To bring up a more specific example, in [9, Lemma 12] it is stated that the CCZ class of $f(x)=x^{3}$ on $\mathbb{F}_{2^{5}}$ contains three EA classes, and this is based on the classification of 64 WZ spaces that according to [9] were found experimentally. Here we can give a computer-free description of these spaces: 32 of them are obtained from Proposition 4.3 with $m=n=5$, and the remaining 32 of them are obtained from Proposition 4.1 with $S=\mu \mathbb{F}_{2}$ where $\mu \in \mathbb{F}_{2^{5}}$.

\subsection*{6.2 Construction of new APN permutations}
If we know two trivially intersecting WZ spaces for an APN function $f$ : $\mathbb{F}_{2^{n}} \rightarrow \mathbb{F}_{2^{n}}$, then Proposition 2.2 allows us to construct an APN permutation $f^{\prime}$ of $\mathbb{F}_{2^{n}}$. Then $f^{\prime}$ is CCZ equivalent to $f$, but in general it need not be EA equivalent to it.

Just for illustration we present a simple numerical example. For $n=$ 9, the algebraic degree of $f(x)=x^{3}$ is 2 , and the algebraic degree of its compositional inverse $g(x)=x^{1 / 3}$, which is also an APN permutation, is 5 . By applying Proposition 2.2 along with constructions of trivially intersecting WZ spaces provided in Section 5 above, we found APN permutations of $\mathbb{F}_{2}{ }^{9}$\\
of algebraic degrees 2, 4 and 5. Since the algebraic degree is preserved by EA equivalence, the APN permutations of degree 4 are not EA equivalent to $f$ or $g$.

The constructions in Section 5 work in arbitrary odd dimensions and for all Gold APN functions. It will be interesting to investigate how many EA inequivalent APN permutations they provide.

\section*{Acknowledgement}
We thank Claude Carlet for pointing to us the current form of Proposition 3.4 which was previously stated in a more special form. We also thank Léo Perrin for releasing an updated version of the sboxU software package [15] that is more widely compatible with various computer operating systems.

\section*{References}
[1] C. Beierle, C. Carlet, G. Leander, L. Perrin, A further study of quadratic APN permutations in dimension nine. arXiv preprint arXiv:2104.08008, 2021.\\[0pt]
[2] C. Beierle, G. Leander, New instances of quadratic APN functions. CoRR, abs/2009.07204, 2020.\\[0pt]
[3] X. Bonnetain, L. Perrin, S. Tian, Anomalies and vector space search: tools for S-box analysis. In: S. Galbraith, S. Moriai (eds) Advances in Cryptology - ASIACRYPT 2019. Lecture Notes in Computer Science, vol. 11921, pp. 196-223. Springer, Cham, 2019.\\[0pt]
[4] K.A. Browning, J.F. Dillon, R.E. Kibler, and M.T. McQuistan, APN polynomials and related codes. Journal of Combinatorics, Information and System Science, Special Issue in honor of Prof. D.K Ray-Chaudhuri on the occasion of his 75 th birthday, vol. 34, no. 1-4, pp. 135-159, (2009).\\[0pt]
[5] K.A. Browning, J.F. Dillon, M.T. McQuistan, A.J. Wolfe, An APN permutation in dimension six. Finite fields: theory and applications, 33-42, Contemp. Math., 518, Amer. Math. Soc., Providence, RI, 2010.\\[0pt]
[6] L. Budaghyan, M. Calderini, C. Carlet, On known constructions of APN and AB functions and their relation to each other. IACR Cryptol. ePrint Arch., 2020:1444, 2020.\\[0pt]
[7] L. Budaghyan, M. Calderini, I. Villa, On equivalence between known families of quadratic APN functions. Finite Fields and Their Applications 66 (2020), 101704.\\[0pt]
[8] L. Budaghyan, C. Carlet, G. Leander, Two classes of quadratic APN binomials inequivalent to power functions. IEEE Trans. Inf. Theory, 54(9):4218-4229, 2008.\\[0pt]
[9] A. Canteaut, L. Perrin, On CCZ-equivalence, extended-affine equivalence, and function twisting. Finite Fields Appl. 56 (2019), 209-246.\\[0pt]
[10] C. Carlet, Open questions on nonlinearity and on APN functions. Arithmetic of finite fields, 83-107, Lecture Notes in Comput. Sci., 9061, Springer, Cham, 2015.\\[0pt]
[11] C. Carlet, Boolean Functions for Cryptography and Coding Theory. Cambridge University Press, 2020.\\[0pt]
[12] C. Carlet, P. Charpin, V. Zinoviev, Codes, bent functions and permutations suitable for DES-like cryptosystems. Des. Codes Cryptogr. 15 (1998), 125-156.\\[0pt]
[13] H. Dobbertin, Almost perfect nonlinear power functions on $\operatorname{GF}\left(2^{n}\right)$ : A new case for $n$ divisible by 5 . International Conference on Finite Fields and Applications, pages 113-121, 2001.\\[0pt]
[14] R. Gold, Maximal recursive sequences with 3 -valued recursive crosscorrelation functions. IEEE Transactions on Information Theory, 14(1):154-156, 1968.\\[0pt]
[15] M. Joly, L. Perrin, SboxU (software library). \href{https://github.com/lppcrypto/sboxU}{https://github.com/lppcrypto/sboxU} (accessed 19 April 2021)\\[0pt]
[16] J. Lahtonen, G. McGuire, H.N. Ward, Gold and Kasami-Welch functions, quadratic forms, and bent functions. Adv. Math. Commun. 1 (2007), no. 2, 243-250.\\[0pt]
[17] D.R. Stinson, M.B. Paterson, Cryptography. Theory and Practice. Fourth edition. CRC Press, Boca Raton, FL, 2019.\\[0pt]
[18] Y. Yu, M. Wang, Y. Li, A matrix approach for constructing quadratic APN functions. Designs, Codes and Cryptography, vol. 73, no. 2, pp. 587-600, 2014.


\end{document}